%==============================================================================
% REWRITTEN ABSTRACT AND CONCLUSION
% Paper: "Liquidity, Activism Disclosure, and Takeover Premia" — Austin Li
% Date: 2026-02-09
%==============================================================================


%==============================================================================
% ABSTRACT
%==============================================================================

\begin{abstract}
\noindent How does secondary-market liquidity shape corporate control when an
informed blockholder can choose between exit, quiet engagement, and public
activism? I develop a takeover model in which a blockholder observes a private
signal about firm value and selects among four actions: Exit, Passive Hold,
engagement below disclosure thresholds (\emph{Quiet Voice}), or stake-building
with mandatory disclosure (\emph{Public Voice}). A competitive market maker
sets prices from order flow and the disclosure flag; a potential bidder
conditions entry on these market signals.

In equilibrium, prices embed an activism premium---either directly observed
through disclosure or inferred from order flow. I show that liquidity has a
\emph{non-monotone} effect on expected minority takeover gains (under
sufficient boundary conditions, verified in the calibration). Rising
noise-trading intensity simultaneously lowers adverse-selection costs and
erodes the informativeness of order flow, pushing bid incidence and engagement
incentives in opposing directions. The result is a hump-shaped relationship:
moderate liquidity maximizes minority gains, while high liquidity destroys the
inference channel that sustains the activism premium. Threshold-based disclosure
attenuates this channel by shifting engagement from inferred to directly
observable, implying that liquidity regulation is implicitly governance policy.
\end{abstract}


%==============================================================================
% CONCLUSION
%==============================================================================

\section{Conclusion}
\label{sec:conclusion}
%==============================================================================

I develop a model in which secondary-market liquidity, activism disclosure, and
takeover premia are jointly determined in equilibrium. An informed blockholder
chooses among Exit, \emph{Quiet Voice}, and \emph{Public Voice}, while a
competitive market maker prices order flow and a potential bidder decides
whether to enter. The central mechanism is an activism premium embedded in
prices---either directly observed through mandatory disclosure or inferred from
order flow by rational agents. I show that liquidity has a non-monotone effect
on expected minority takeover gains: rising noise-trading intensity lowers
adverse-selection costs but simultaneously erodes the informativeness of order
flow, pushing bid incidence and engagement incentives in opposing directions.
The baseline calibration produces a hump-shaped relationship, with moderate
liquidity maximizing minority gains.

These results carry direct policy implications. Cross-country variation in
disclosure regimes---such as the United Kingdom's lower ownership thresholds
relative to those in the United States---provides natural laboratories for
testing the model's predictions. Lowering the disclosure threshold increases
the share of activism that is directly observable, attenuating the inference
channel through which liquidity shapes takeover premia. Yet lower thresholds
also reduce the returns to quiet engagement, potentially discouraging activism
altogether; the optimal threshold therefore trades off transparency against
incentives. More broadly, policies that affect liquidity---tick sizes, trading
halts, dark pool regulation---alter the noise-trading environment and hence the
way engagement is priced. In markets where activist investors play an important
governance role, liquidity policy is implicitly governance policy. Similarly,
minority shareholder protection depends on equilibrium activism: reforms that
appear protective in partial equilibrium may reduce takeover premia in general
equilibrium by discouraging the engagement that generates them.

Extending the framework to dynamic settings with multiple trading rounds would
capture the gradual position-building that characterizes real activist
campaigns. Introducing multiple blockholders or wolf-pack coordination would
illuminate how collective action reshapes the inference channel. Endogenizing
the blockholder's information acquisition would connect the governance
mechanism to the broader literature on informed trading and price efficiency.
Each extension preserves the core insight that liquidity, disclosure, and
governance are jointly determined---and each offers a richer laboratory for
quantifying the tradeoffs at the heart of market design.
